

\begin{derivation}[从\eqnrefs{eq:bubble-imag-phase-space-dp8}到\eqnrefs{eq:two-body-threshold-beta-dp8}]

正文\eqnrefs{eq:bubble-imag-phase-space-dp8} 表明虚部只有在两个内部粒子同时满足正能质量壳和能量守恒时才存在。对等质量情形进入质心系，能量 $\delta$ 函数固定唯一的径向动量 $k_*$；解出这个根并把它写成相对论速度因子，就得到正文\eqnrefs{eq:two-body-threshold-beta-dp8}。

在质心系 $\mathbf P=0$ 且 $m_1=m_2=m$ 时，
\begin{equation*}
 E_1=E_2=E_k=\sqrt{m^2c^4+c^2k^2},
 \qquad
 \sqrt{s}\,c=2E_k.
\end{equation*}
平方后有
\begin{align*}
 sc^2
 &=4m^2c^4+4c^2k_*^2,\\
 k_*^2
 &=\frac{s-4m^2c^2}{4}.
\end{align*}
因此只有 $s\ge4m^2c^2$ 时 $k_*$ 才为实数。再用 $E_{\rm cm}=\sqrt{s}\,c$，
\begin{align*}
 \beta_*(s)
 &=\frac{ck_*}{E_{\rm cm}/2}
 =\frac{2ck_*}{\sqrt{s}\,c}\\
 &=\sqrt{1-\frac{4m^2c^2}{s}},
\end{align*}
即正文\eqnrefs{eq:two-body-threshold-beta-dp8}。为把阈值因子的整体归一化也固定下来，定义正文\eqnrefs{eq:bubble-imag-phase-space-dp8} 中的径向相空间核
\begin{equation*}
 \rho_2(s)
 \equiv
 \int\frac{\dd^3\mathbf k}{(2\pi\hbar)^3}
 \frac{c^2}{4E_k^2}
 \delta\!\left(\sqrt{s}\,c-2E_k\right).
\end{equation*}
利用
\begin{equation*}
 \left.\frac{\dd(2E_k)}{\dd k}\right|_{k_*}
 =\frac{2c^2k_*}{E_*},
 \qquad E_*\equiv E_{k_*}=\frac{\sqrt{s}\,c}{2},
\end{equation*}
径向 $\delta$ 积分严格给出
\begin{align*}
 \rho_2(s)
 &=\frac{4\pi k_*^2}{(2\pi\hbar)^3}
 \frac{c^2}{4E_*^2}
 \frac{E_*}{2c^2k_*}\\
 &=\frac{k_*}{16\pi^2\hbar^3E_*}
 =\frac{\beta_*(s)}{16\pi^2\hbar^3c}.
\end{align*}
因此连续谱在阈值处按确定的线性因子 $\beta_*(s)$ 从零打开。
\end{derivation}













\begin{derivation}[从\eqnrefs{eq:ssb-complex-scalar-lag-dp11}到\eqnrefs{eq:ssb-higgs-goldstone-masses-dp11}]

正文\eqnrefs{eq:ssb-complex-scalar-lag-dp11} 给出了只依赖 \(|\varphi_{\rm ssb}|\) 的对称势。先由驻点条件确定真空半径，再在所选真空附近展开到二次阶；径向与切向二次系数的差别最终给出正文\eqnrefs{eq:ssb-higgs-goldstone-masses-dp11} 的有质量径向模和无质量切向模。


令 \(r^2\equiv\varphi_{\rm ssb}^*\varphi_{\rm ssb}\)，则
\begin{equation*}
 V(r)=-\mu_{\rm ssb}^2r^2+\frac\lambda{\hbar c}r^4.
\end{equation*}
极值条件为
\begin{align*}
 \frac{\dd V}{\dd r}
 &=-2\mu_{\rm ssb}^2r+\frac{4\lambda}{\hbar c}r^3,\\
 &=2r\left(-\mu_{\rm ssb}^2+\frac{2\lambda}{\hbar c}r^2\right)=0.
\end{align*}
除去不稳定的 \(r=0\)，稳定极小值满足
\begin{equation*}
 r^2=\frac{\mu_{\rm ssb}^2\hbar c}{2\lambda}.
\end{equation*}
定义 \(r=v/\sqrt2\) 即得到
\begin{equation*}
 v^2=\frac{\mu_{\rm ssb}^2\hbar c}{\lambda}.
\end{equation*}

将真空附近参数化代入
\begin{equation*}
 \varphi_{\rm ssb}=\frac1{\sqrt2}(v+h+\ii\pi),
 \qquad
 \varphi_{\rm ssb}^*\varphi_{\rm ssb}=\frac12[(v+h)^2+\pi^2].
\end{equation*}
保留到二次阶：
\begin{align*}
 V
 ={}&-\frac{\mu_{\rm ssb}^2}{2}
 (v^2+2vh+h^2+\pi^2)\\
 &+\frac\lambda{4\hbar c}
 (v^2+2vh+h^2+\pi^2)^2.
\end{align*}
线性 \(h\) 项系数为
\begin{equation*}
 -\mu_{\rm ssb}^2v+\frac{\lambda v^3}{\hbar c}=0,
\end{equation*}
二次项为
\begin{equation*}
 V^{(2)}=\mu_{\rm ssb}^2h^2+0\times\pi^2.
\end{equation*}
与 \(\tfrac12(m^2c^2/\hbar^2)\phi^2\) 比较即得\eqnrefs{eq:ssb-higgs-goldstone-masses-dp11}。
\end{derivation}

% 守恒局域荷的无能隙判据：已迁入 sections/03qft-ssb-higgs-general-dp11.tex 的结论框及独立证明。


\begin{derivation}[从\eqnrefs{eq:abelian-higgs-polar-dp11}到\eqnrefs{eq:abelian-higgs-vector-mass-dp11}]
从
\[
\varphi_{\rm ssb}(x)
=\frac{v+h(x)}{\sqrt2}\ee^{\ii\pi(x)/v},
\qquad
D_\mu=\partial_\mu+\ii g_{\rm G}\mathscr C_\mu
\]
出发。局域 $U(1)$ 变换
\[
\varphi_{\rm ssb}\to \ee^{\ii\alpha(x)}\varphi_{\rm ssb},
\qquad
\mathscr C_\mu\to
\mathscr C_\mu-\frac1{g_{\rm G}}\partial_\mu\alpha
\]
要求极坐标变量变成
\[
h\to h,
\qquad
\pi\to\pi+v\alpha.
\]
因此组合
\[
\mathcal B_\mu
\equiv
\mathscr C_\mu+\frac1{g_{\rm G}v}\partial_\mu\pi
\]
满足 $\mathcal B_\mu\to\mathcal B_\mu$，它才是破缺相位中的规范不变矢量变量。

显式求协变导数：
\begin{align*}
D_\mu\varphi_{\rm ssb}
&=\frac{\ee^{\ii\pi/v}}{\sqrt2}
\left[
\partial_\mu h
+\ii(v+h)
\left(\frac{\partial_\mu\pi}{v}+g_{\rm G}\mathscr C_\mu\right)
\right]\\
&=\frac{\ee^{\ii\pi/v}}{\sqrt2}
\left[
\partial_\mu h
+\ii g_{\rm G}(v+h)\mathcal B_\mu
\right].
\end{align*}
取模平方时，实的径向项与纯虚的相位项交叉项严格相消：
\[
(D_\mu\varphi)^*(D^\mu\varphi)
=\frac12(\partial h)^2
+\frac{g_{\rm G}^2}{2}(v+h)^2\mathcal B_\mu\mathcal B^\mu.
\]
展开 $(v+h)^2=v^2+2vh+h^2$。与矢量二次型有关的纯二次项是
\[
\frac{g_{\rm G}^2v^2}{2}\mathcal B_\mu\mathcal B^\mu.
\]
在上述 SI 归一化下，Proca 质量项写成
\[
\frac{\hbar c}{2}
\left(\frac{m_Ac}{\hbar}\right)^2
\mathcal B_\mu\mathcal B^\mu.
\]
两者是同一个二次算符，所以系数必须相等：
\[
g_{\rm G}^2v^2
=\hbar c\left(\frac{m_Ac}{\hbar}\right)^2,
\]
从而
\[
\boxed{m_Ac^2=g_{\rm G}\sqrt{\hbar c}\,v.}
\]
Goldstone 场 $\pi$ 通过 $\mathcal B_\mu$ 重组为有质量矢量场的纵向偏振。破缺前自由度计数为“复标量的两个实自由度 + 无质量矢量的两个横向偏振”共四个；破缺后为“径向标量 $h$ 一个 + 有质量矢量 $\mathcal B_\mu$ 三个偏振”仍共四个。这个计数是上述变量重组没有丢失物理自由度的独立检查。
\end{derivation}

\begin{derivation}[从\eqnrefs{eq:gauge-goldstone-mixing-dp11}到\eqnrefs{eq:rxi-goldstone-mass-abelian-dp42}]
Cartesian 展开给出
\[
\Lag_{\rm mix}=g_{\rm G}v\,\mathscr C^\mu\partial_\mu\pi,
\qquad
F=\partial\!\cdot\!\mathscr C-\xi\frac{g_{\rm G}v}{\hbar c}\pi .
\]
规范固定项展开为
\[
\Lag_{\rm gf}
=-\frac{\hbar c}{2\xi}(\partial\!\cdot\!\mathscr C)^2
+g_{\rm G}v\,\pi\,\partial\!\cdot\!\mathscr C
-\frac{\xi g_{\rm G}^2v^2}{2\hbar c}\pi^2 .
\]
对中间项分部积分，
\[
g_{\rm G}v\,\pi\,\partial_\mu\mathscr C^\mu
=-g_{\rm G}v\,\mathscr C^\mu\partial_\mu\pi+\text{边界项},
\]
故它与 $\Lag_{\rm mix}$ 相消。剩余 $\pi^2$ 项与标准质量项比较：
\[
-\frac{\xi g_{\rm G}^2v^2}{2\hbar c}\pi^2
=-\frac12\left(\frac{m_\pi c}{\hbar}\right)^2\pi^2,
\qquad
\left(\frac{m_Ac}{\hbar}\right)^2=\frac{g_{\rm G}^2v^2}{\hbar c},
\]
于是
\[
\boxed{m_\pi^2=\xi m_A^2}.
\]
$\pi$ 的质量依赖规范参数 $\xi$，而物理规范玻色子质量 $m_A$ 不随 $\xi$ 改变。
\end{derivation}

\begin{derivation}[从\eqnrefs{eq:all-incoming-amplitude-dp9}到\eqnrefs{eq:physical-from-all-incoming-dp9}]
全入射约定把每一条外腿的四动量都定义成指向相互作用区。设全入射核的四个参数为 $k_i$，则正文\eqnrefs{eq:all-incoming-amplitude-dp9} 要求
\[
k_1+k_2+k_3+k_4=0.
\]
对物理过程 $1+2\to3+4$，实验四动量满足
\[
p_1+p_2=p_3+p_4.
\]
前两条腿本来就流入，所以
\[
k_1=p_1,\qquad k_2=p_2.
\]
末态 3、4 的物理箭头流出相互作用区；若仍要求 $k_i$ 全部向内，必须反号：
\[
k_3=-p_3,\qquad k_4=-p_4.
\]
于是全入射的动量守恒自动变成
\begin{align*}
k_1+k_2+k_3+k_4
&=p_1+p_2-p_3-p_4\\
&=0,
\end{align*}
与物理守恒完全一致。将四个 $k_i$ 代回同一个截肢四点核，便得到正文\eqnrefs{eq:physical-from-all-incoming-dp9}：
\[
\mathcal M_{12\to34}
=\mathcal A(p_1,p_2,-p_3,-p_4).
\]
\end{derivation}



\begin{derivation}[从\eqnrefs{eq:2to2-cross-section-master-si-r9}到\eqnrefs{eq:relativistic-to-nr-potential-matching-dp75}]
低速弹性区有 $|\mathbf p_f|=|\mathbf p_i|$。定义非相对论振幅为阈值极限
\begin{equation*}
\mathcal M_{\rm rel}^{\rm NR}(\mathbf q)
\equiv\lim_{|\mathbf p_i|/(m_{A,B}c)\to0}\mathcal M_{\rm rel},
\qquad
\lim_{|\mathbf p_i|\to0}\sqrt{s}=(m_A+m_B)c.
\end{equation*}
因此对截面主式取同一极限得到
\begin{equation*}
 \frac{\dd\sigma}{\dd\Omega}
 =\frac{\hbar^2}{64\pi^2(m_A+m_B)^2c^2}
 |\mathcal M_{\rm rel}^{\rm NR}(\mathbf q)|^2.
\end{equation*}
另一方面，对\eqnrefs{eq:sch-lippmann-r13} 作一次 Born 迭代，采用同一出射波相位约定可写
\begin{equation*}
 f_B(\mathbf q)
 =-\frac{m_{\rm red}}{2\pi\hbar^2}
 \widetilde V_{\rm NR}(\mathbf q),
 \qquad
 \frac{\dd\sigma}{\dd\Omega}=|f_B|^2,
\end{equation*}
其中 $m_{\rm red}=m_Am_B/(m_A+m_B)$。与相对论截面对应的散射振幅归一化为
\begin{equation*}
 f_{\rm rel}(\mathbf q)
 =\frac{\hbar}{8\pi(m_A+m_B)c}
 \mathcal M_{\rm rel}^{\rm NR}(\mathbf q),
\end{equation*}
其相位由前面固定的 $S$ 矩阵约定与出射 Green 核共同确定。令 $f_B=f_{\rm rel}$，得到
\begin{align*}
 \widetilde V_{\rm NR}(\mathbf q)
 &=-\frac{\hbar^3}{4m_{\rm red}(m_A+m_B)c}
 \mathcal M_{\rm rel}^{\rm NR}(\mathbf q)\\
 &=-\frac{\hbar^3}{4m_Am_Bc}
 \mathcal M_{\rm rel}^{\rm NR}(\mathbf q),
\end{align*}
即\eqnrefs{eq:relativistic-to-nr-potential-matching-dp75}。
\end{derivation}

% 局域反项的递归减除：已迁入 sections/03qft-renorm-allorders-ope-dp11.tex 的结论框及独立证明。






















\begin{derivation}[从\eqnrefs{eq:phi4-bubble-integral-r68}到\eqnrefs{eq:phi4-bubble-result-dp10}]
\eqnrefs{eq:dim-reg-measure-general-dp10} 中显式 \(\hbar\) 的消去可由变量替换直接看出。若从逆长度圈变量 \(\kappa^\mu\) 开始，定义物理动量

\begin{equation*}
 \ell^\mu=\hbar\kappa^\mu,
 \qquad
 \mu_{\rm DR}=\hbar\mu_\kappa,
\end{equation*}
则
\begin{equation*}
 \dd^d \kappa=\frac{\dd^d \ell}{\hbar^d},
 \qquad
 (\kappa^2-M^2)^{-2}
 =\frac{\hbar^4}{(\ell^2-m_R^2c^2)^2},
\end{equation*}
其中 \(M=m_Rc/\hbar\)。维数正则化尺度给
\begin{equation*}
 \mu_\kappa^{2\epsilon}
 =\frac{\mu_{\rm DR}^{2\epsilon}}{\hbar^{2\epsilon}}.
\end{equation*}
由于 \(d=4-2\epsilon\)，总的 \(\hbar\) 次数为
\begin{equation*}
 4-d-2\epsilon=0.
\end{equation*}
因此变量换元后的物理四动量形式为
\begin{equation*}
 \mu_{\rm DR}^{2\epsilon}
 \int\frac{\dd^d \ell}{(2\pi)^d}
 \frac{\ii}{(\ell^2-m_R^2c^2+\ii0^+)
 ((\ell+p)^2-m_R^2c^2+\ii0^+)},
\end{equation*}
其中显式 $\hbar$ 已由测度、分母和正则化尺度的幂次精确抵消。

用 Feynman 参数
\begin{equation*}
 \frac1{AB}=\int_0^1\frac{\dd x}{[xA+(1-x)B]^2}
\end{equation*}
并令 \(L=\ell+xp\)，分母化为
\begin{equation*}
 L^2-\Delta_s(x)+\ii0^+,
 \qquad
 \Delta_s(x)=m_R^2c^2-x(1-x)s.
\end{equation*}
于是
\begin{equation*}
 I(s)=\int_0^1\dd x\,
 \mu_{\rm DR}^{2\epsilon}
 \int\frac{\dd^d L}{(2\pi)^d}
 \frac{\ii}{(L^2-\Delta_s+\ii0^+)^2}.
\end{equation*}
做 Wick 转动 \(L^0=\ii L_E^0\)：
\begin{equation*}
 \dd^d L=\ii\dd^d L_E,
 \qquad
 L^2-\Delta+\ii0^+=-(L_E^2+\Delta-\ii0^+),
\end{equation*}
故
\begin{equation*}
 I(s)=-\int_0^1\dd x\,
 \mu_{\rm DR}^{2\epsilon}
 \int\frac{\dd^d L_E}{(2\pi)^d}
 \frac1{(L_E^2+\Delta_s-\ii0^+)^2}.
\end{equation*}
利用欧氏母积分
\begin{equation*}
 \int\frac{\dd^d L_E}{(2\pi)^d}
 \frac1{(L_E^2+\Delta)^2}
 =\frac1{(4\pi)^{d/2}}
 \Gamma\!\left(2-\frac d2\right)
 \Delta^{d/2-2},
\end{equation*}
代入 \(d=4-2\epsilon\)：
\begin{equation*}
 I(s)=-\int_0^1\dd x\,
 \frac1{(4\pi)^2}
 \Gamma(\epsilon)(4\pi)^\epsilon
 \left(\frac{\mu_{\rm DR}^2}{\Delta_s-\ii0^+}\right)^\epsilon.
\end{equation*}
逐项展开
\begin{align*}
 \Gamma(\epsilon)&=\frac1\epsilon-\gamma_E+\order(\epsilon),\\
 (4\pi)^\epsilon&=1+\epsilon\ln4\pi+\order(\epsilon^2),\\
 \left(\frac{\mu_{\rm DR}^2}{\Delta_s-\ii0^+}\right)^\epsilon
 &=1+\epsilon\ln\frac{\mu_{\rm DR}^2}{\Delta_s-\ii0^+}
 +\order(\epsilon^2).
\end{align*}
相乘到有限阶并执行参数积分，得到
\begin{equation*}
 I(s)=-\frac1{16\pi^2}
 \left[
 \frac1{\bar\epsilon}
 -\int_0^1\dd x\,
 \ln\frac{\Delta_s(x)-\ii0^+}{\mu_{\rm DR}^2}
 \right]+
 \order(\epsilon),
\end{equation*}
即\eqnrefs{eq:phi4-bubble-result-dp10}。当 \(\Delta_s(x)\) 穿过零时，\(-\ii0^+\) 决定对数的虚部，并与在壳中间态的 Cutkosky 不连续度一致。
\end{derivation}

\begin{derivation}[从\eqnrefs{eq:phi4-bare-coupling-r68}到\eqnrefs{eq:phi4-beta-one-loop-dp10}]
记

\begin{equation*}
 a_1\equiv\frac{3}{32\pi^2},
 \qquad
 P_\epsilon\equiv\frac1{\bar\epsilon}.
\end{equation*}
一圈阶裸耦合为
\begin{equation*}
 g_{4,B}
 =\mu^{2\epsilon}
 \left[g_4(\mu)+a_1P_\epsilon g_4^2(\mu)
 +\order(g_4^3)\right].
\end{equation*}
按正文\eqnrefs{eq:phi4-beta-one-loop-dp10} 中 $\beta_{g_4}$ 的定义，$d=4-2\epsilon$ 时固定裸耦合的总尺度导数写成
\begin{equation*}
 \mu\frac{\dd g_4}{\dd\mu}\Big|_{g_{4,B}}
 =-2\epsilon g_4+\beta_{g_4}(g_4).
\end{equation*}
对裸表达式求导时，反项同样通过 \(g_4(\mu)\) 依赖于 \(\mu\)：
\begin{align*}
 0={}&2\epsilon
 \left(g_4+a_1P_\epsilon g_4^2\right)\\
 &+\left(1+2a_1P_\epsilon g_4\right)
 \left(-2\epsilon g_4+\beta_{g_4}\right)
 +\order(g_4^3).
\end{align*}
保留到 \(g_4^2\) 后，\(2\epsilon g_4\) 项相消，并得到
\begin{equation*}
 0=\beta_{g_4}
 -2a_1\epsilon P_\epsilon g_4^2
 +\order(g_4^3,\epsilon g_4^2).
\end{equation*}
因为
\begin{equation*}
 \epsilon P_\epsilon
 =\epsilon\left(\frac1\epsilon-\gamma_E+\ln4\pi\right)
 \longrightarrow1,
\end{equation*}
四维极限给出
\begin{equation*}
 \beta_{g_4}=2a_1g_4^2+
\order(g_4^3)
 =\frac{3g_4^2}{16\pi^2}+
\order(g_4^3).
\end{equation*}
这里关键的因子 2 来自 \(d=4-2\epsilon\) 下的 \(\mu^{2\epsilon}\)；若把反项对 \(g_4(\mu)\) 的隐含尺度依赖漏掉，就会错误地得到相反符号。
\end{derivation}











% 定态散射的能量壳关系：已迁入 sections/02a0-exact-scattering.tex 的结论框及独立证明。






\begin{derivation}[从\eqnrefs{eq:phi4-first-dyson-r68}到\eqnrefs{eq:phi4-tree-green-r3}]

正文\eqnrefs{eq:phi4-first-dyson-r68} 是 $\varphi^4$ 理论四点函数的一阶 Dyson 插入。连通树级项要求每个外场都与顶点 $z$ 上四个场中的一个且仅一个收缩，因此共有 $4!$ 种相同配对；它们与作用量中的 $1/4!$ 抵消后得到正文\eqnrefs{eq:phi4-tree-green-r3}。



四点函数的一阶 Dyson 项为
\begin{equation*}
 G^{(4,1)}(x_1,x_2,x_3,x_4)
 =-\frac{\ii g_4}{4!\,\hbar^2c^2}
 \int\dd^4 z\,
 \langle0|\Tord\{\varphi_1\varphi_2\varphi_3\varphi_4\varphi^4(z)\}|0\rangle .
\end{equation*}
要得到“一个顶点连接四个外点”的连通树图，每个外部场必须与 $z$ 点四个场中的一个且仅一个收缩。先给 $\varphi_1$ 选 $z$ 点场有 $4$ 种；再给 $\varphi_2$ 选剩余场有 $3$ 种；之后分别有 $2$ 种和 $1$ 种，所以
\begin{equation*}
 N_{\rm contr}=4\times3\times2\times1=4!.
\end{equation*}
每一种收缩都给出同一个核乘积
\begin{equation*}
 \Delta_F(x_1-z)\Delta_F(x_2-z)\Delta_F(x_3-z)\Delta_F(x_4-z).
\end{equation*}
因此 $4!$ 个相同项与相互作用中的 $1/4!$ 抵消，得到
\begin{equation*}
 G_c^{(4)}
 =-\frac{\ii g_4}{\hbar^2c^2}
 \int\dd^4 z\,\prod_{a=1}^4\Delta_F(x_a-z),
\end{equation*}
即正文\eqnrefs{eq:phi4-tree-green-r3}。

再把四条外线分别 Fourier 展开。四个传播子的指数中所有 $z$ 依赖合成
\begin{equation*}
 \exp\!\left[-\frac{\ii}{\hbar}
 (p_1+p_2-p_3-p_4)\cdot z\right].
\end{equation*}
于是顶角位置积分给
\begin{equation*}
 \int\dd^4 z\,
 \ee^{-\ii(p_1+p_2-p_3-p_4)\cdot z/\hbar}
 =(2\pi\hbar)^4\delta^{(4)}(p_1+p_2-p_3-p_4).
\end{equation*}
所以顶角四动量守恒正是局域相互作用对顶角位置积分后的 Fourier $\delta$ 函数。
\end{derivation}

\begin{derivation}[从\eqnrefs{eq:golden-sinc-r3}到\eqnrefs{eq:fermi-golden-r3}]
记 $\Delta E=E_f-E_i$。一阶振幅为

\begin{align*}
 c_f^{(1)}(T)
 &=-\frac{\ii}{\hbar}V_{fi}
 \int_{-T/2}^{T/2}\dd t\,\ee^{\ii\Delta E t/\hbar}\\
 &=-\frac{\ii}{\hbar}V_{fi}
 \left[\frac{\hbar}{\ii\Delta E}
 \ee^{\ii\Delta E t/\hbar}\right]_{-T/2}^{T/2}\\
 &=-\frac{2\ii V_{fi}}{\Delta E}
 \sin\!\left(\frac{\Delta E T}{2\hbar}\right).
\end{align*}
把 $\sinc x\equiv\sin x/x$ 提出，得到
\begin{equation*}
 |c_f^{(1)}(T)|^2
 =\frac{|V_{fi}|^2T^2}{\hbar^2}
 \sinc^2\!\left(\frac{\Delta E T}{2\hbar}\right),
\end{equation*}
即\eqnrefs{eq:golden-sinc-r3}。它在有限 $T$ 时不是 $\delta$ 函数；主峰宽度量级是 $\hbar/T$。

对连续末态求和，概率为
\begin{equation*}
 P_{i\to f}(T)
 =\int\dd E\,\nu_f(E)
 \frac{|V_{fi}(E)|^2T^2}{\hbar^2}
 \sinc^2\!\left[\frac{(E-E_i)T}{2\hbar}\right].
\end{equation*}
这里 $\nu_f(E)$ 是末态每单位能量的态密度，SI 量纲为 $\mathrm{J^{-1}}$。有限时间谱核与 Dirac $\delta$ 分布的对应由下列核固定：
\begin{equation*}
 K_T(\Delta E)
 \equiv
 \frac{T}{2\pi\hbar}
 \sinc^2\!\left(\frac{\Delta E\,T}{2\hbar}\right).
\end{equation*}
它具有能量倒数量纲。作变量代换 $x=\Delta E T/(2\hbar)$ 后，
\begin{align*}
 \int_{-\infty}^{+\infty}\dd(\Delta E)\,K_T(\Delta E)
 &=\frac{T}{2\pi\hbar}
 \frac{2\hbar}{T}
 \int_{-\infty}^{+\infty}\dd x\,\sinc^2x\\
 &=\frac1\pi\int_{-\infty}^{+\infty}\dd x\,\sinc^2x=1.
\end{align*}
因此 $K_T$ 是面积固定为一、宽度随 $T^{-1}$ 收缩的归一化谱核；对任意在 $E_i$ 附近平滑的测试函数 $F(E)$，$T\to\infty$ 时满足
\begin{equation*}
 \int\dd E\,F(E)K_T(E-E_i)\longrightarrow F(E_i),
\end{equation*}
也就是 $K_T(\Delta E)\to\delta(\Delta E)$ 的分布极限。
令
\begin{equation*}
F(E)\equiv |V_{fi}(E)|^2\nu_f(E).
\end{equation*}
由 $K_T$ 的定义，有限时间概率可以严格改写为
\begin{equation*}
P_{i\to f}(T)
=\frac{2\pi T}{\hbar}
\int\dd E\,F(E)K_T(E-E_i).
\end{equation*}
再定义有限时间余项
\begin{equation*}
\delta F_T
\equiv
\int\dd E\,[F(E)-F(E_i)]K_T(E-E_i),
\end{equation*}
利用 $\int K_T\dd E=1$ 得到精确等式
\begin{equation*}
P_{i\to f}(T)
=\frac{2\pi T}{\hbar}[F(E_i)+\delta F_T].
\end{equation*}
前面已经证明 $K_T\to\delta$ 的分布极限，所以对满足该近似恒等核收敛条件的连续 $F$，有 $\lim_{T\to\infty}\delta F_T=0$。因此长时间跃迁率严格由极限定义
\begin{equation*}
\lim_{T\to\infty}\frac{P_{i\to f}(T)}{T}
=\frac{2\pi}{\hbar}F(E_i).
\end{equation*}
也就是
\begin{equation*}
 w_{i\to f}=\frac{2\pi}{\hbar}|V_{fi}|^2\nu_f(E_i),
\end{equation*}
即 Fermi 黄金律~\eqnrefs{eq:fermi-golden-r3}。能量 $\delta$ 分布由有限时间谱核与连续谱卷积的长时间极限产生；有限 $T$ 时应保留具有有限宽度的 $\sinc^2$ 核。
\end{derivation}





\begin{derivation}[从\eqnrefs{eq:lsz-pole-plus-cont-r7}到\eqnrefs{eq:phi4-lsz-result-r7}]
先只处理一条外腿，因为多腿公式只是把同一个极点提取操作重复 $n$ 次。设 $\zeta=p^2-m_{\rm phys}^2c^2$，并把 $\zeta$ 暂时视为复变量。对时序关联函数中与坐标 $x$ 对应的那一个场插入完整态，稳定单粒子态给出
\begin{equation*}
 \langle\Omega|\varphi_H(0)|p,\lambda\rangle
 =Z^{1/2}\,\mathcal N_p,
\end{equation*}
其中 $\mathcal N_p$ 是由全书协变态归一化固定的运动学因子。把该单粒子贡献与其余所有中间态分开，Fourier 变换后必然具有形式
\begin{equation*}
 \widetilde G_{n,c}(p,\ldots)
 =\frac{\ii\hbar^3c\,Z^{1/2}}
 {\zeta}\,
 \widetilde F_{n-1}(p;\ldots)
 +\widetilde R^{(1)}_n(p,\ldots),
\end{equation*}
这里 $\widetilde F_{n-1}$ 是已经把这一条外腿换成归一化单粒子态后的矩阵元；$\widetilde R^{(1)}_n$ 包含连续谱和在这一外部通道中正则的项。稳定粒子极点与多粒子阈值分离，所以在 $\zeta=0$ 的一个小邻域内
\begin{equation*}
 \lim_{\zeta\to0}\zeta\,\widetilde R^{(1)}_n=0.
\end{equation*}
因此极点留数不是一句“外腿传播子可以删掉”，而是如下复分析极限：
\begin{align*}
 \lim_{\zeta\to0}
 Z^{-1/2}\frac{\zeta}{\ii\hbar^3c}
 \widetilde G_{n,c}
 &=\lim_{\zeta\to0}
 \left[
 \widetilde F_{n-1}
 +Z^{-1/2}\frac{\zeta}{\ii\hbar^3c}
 \widetilde R^{(1)}_n
 \right]\\
 &=\widetilde F_{n-1}(p_{\rm on};\ldots).
\end{align*}
Feynman 的 $+\ii0^+$ 是实动量边界值的处方；截肢本身应理解为先在复 $\zeta$ 平面取极点留数，再回到质量壳，因此不会出现把 $\zeta/(\zeta+\ii0^+)$ 当作普通数值比值的歧义。

对第二、第三直到第 $n$ 条外腿重复同一操作。若定义 $\zeta_i=p_i^2-m_i^2c^2$，则联合极点邻域中的多变量 Laurent 展开可写成
\begin{align*}
 \widetilde G_{n,c}
 ={}&(2\pi\hbar)^4\delta^{(4)}\!\left(\sum_i\sigma_i^{\rm io}p_i\right)
 \left[\prod_i
 \frac{\ii\hbar^3c\,Z_i^{1/2}}{\zeta_i}\right]
 \ii\mathfrak A_{fi}(\{p_i\})
 +\widetilde R_n,\\
 &\lim_{\zeta_1,\ldots,\zeta_n\to0}
 \left(\prod_i\zeta_i\right)\widetilde R_n=0.
\end{align*}
第二行正是“余项至少缺少一条完整外腿极点”的数学含义。于是
\begin{align*}
 &\lim_{\zeta_1,\ldots,\zeta_n\to0}
 \left[\prod_iZ_i^{-1/2}
 \frac{\zeta_i}{\ii\hbar^3c}\right]
 \widetilde G_{n,c}\\
 &\qquad=(2\pi\hbar)^4\delta^{(4)}\!\left(\sum_i\sigma_i^{\rm io}p_i\right)
 \ii\mathfrak A_{fi}(\{p_{i,\rm on}\}),
\end{align*}
这就得到正文的动量空间 LSZ 公式。注意截肢和在壳极限的顺序来自极点留数的定义：先乘 $\zeta_i$ 提取留数，再令 $\zeta_i\to0$；若反过来，原 Green 函数正好落在极点上。

现在把这一抽象结构用于树级 $\varphi^4$。一次 Dyson 插入给
\begin{equation*}
 -\frac{\ii}{\hbar c}
 \int\dd^4x\,
 \frac{g_4}{4!\,\hbar c}\,
 \varphi_I(x)^4.
\end{equation*}
把四个外部场分别与顶角处四个场收缩有 $4!$ 种双射，恰好消掉相互作用中的 $1/4!$。每条收缩给一个自由 Feynman 传播子，所以 Fourier 变换后的连通四点函数为
\begin{equation*}
 \widetilde G^{(1)}_{4,c}
 =-\frac{\ii g_4}{\hbar^2c^2}
 (2\pi\hbar)^4\delta^{(4)}(p_1+p_2-p_3-p_4)
 \prod_{a=1}^{4}
 \frac{\ii\hbar^3c}{p_a^2-m^2c^2+\ii0^+}.
\end{equation*}
在树级 $Z=1$。令 $\zeta_a=p_a^2-m^2c^2$，逐腿乘逆极点因子后，乘积中的四条外腿逐一消去：
\begin{align*}
 &\left[\prod_{a=1}^{4}
 \frac{\zeta_a}{\ii\hbar^3c}\right]
 \widetilde G^{(1)}_{4,c}\\
 &\quad=-\frac{\ii g_4}{\hbar^2c^2}
 (2\pi\hbar)^4\delta^{(4)}(p_1+p_2-p_3-p_4)
 \prod_{a=1}^{4}
 \frac{\zeta_a}{\zeta_a+\ii0^+}.
\end{align*}
按上面说明的极点留数顺序取极限，四个因子都贡献单位留数，于是得到正文\eqnrefs{eq:phi4-lsz-result-r7}。因此树级顶角规则不是另记的一条口诀：$4!$ 来自 Wick 收缩计数，四条传播子来自外部收缩，而 LSZ 的四个逆极点恰好把这些外腿传播子截去，只留下局域顶角系数。
\end{derivation}

% 稳定标量单腿的约化恒等式：已迁入 detailed/ch03.tex 的结论框及独立证明。






\begin{derivation}[从\eqnrefs{eq:heat-kernel-a4-flat-dp11}到\eqnrefs{eq:heat-kernel-divergence-master-dp11}]

对正定欧几里得二次算符 $\Delta$，一圈实玻色 Gaussian 行列式可用正文已经建立的固有时表示写成
\[
\frac12\Tr\ln\Delta
=-\frac12\mu_{\rm DR}^{2\epsilon}
\int_0^\infty\frac{\dd s}{s}\,
\Tr \ee^{-s\Delta},
\]
其中与场无关的加法常数不影响局域反项。把正文\eqnrefs{eq:heat-kernel-a4-flat-dp11} 代入，并只保留含 $s^2$ 的四维局域系数，得到
\begin{align*}
\left.\frac12\Tr\ln\Delta\right|_{a_4}
={}&-\frac12\mu_{\rm DR}^{2\epsilon}
\int\dd^d x\,
\tr a_4(x)\\
&\times\int_0^\infty\frac{\dd s}{s}
\frac{s^2}{(4\pi s)^{d/2}}.
\end{align*}
取 $d=4-2\epsilon$ 后，短固有时端的幂次变为
\[
\frac1s\frac{s^2}{s^{2-\epsilon}}
=s^{-1+\epsilon}.
\]
因此紫外极点只需从 $s\to0^+$ 端读取。任取固定分界尺度 $s_0$，
\begin{align*}
\mu_{\rm DR}^{2\epsilon}
\int_0^{s_0}\dd s\,s^{-1+\epsilon}
&=\mu_{\rm DR}^{2\epsilon}
\frac{s_0^\epsilon}{\epsilon}\\
&=\frac1\epsilon
+\ln(\mu_{\rm DR}^2s_0)+O(\epsilon).
\end{align*}
所以极点系数与任意分界 $s_0$ 无关，而长固有时部分只改变有限项。代入
\[
 a_4(x)
 =\frac12E^2
 +\frac1{12}\Omega_{\mu\nu}\Omega^{\mu\nu}
 +\frac16\nabla^2E
\]
并使用无边界体积分
\(
\int\dd^4 x\,\nabla^2E=0
\)，得到
\[
\left.\frac12\Tr\ln\Delta\right|_{\rm pole}
=-\frac1{2(4\pi)^2}\frac1\epsilon
\int\dd^4 x\,
\tr\left[
\frac12E^2+\frac1{12}\Omega_{\mu\nu}\Omega^{\mu\nu}
\right].
\]
最后使用正文已经固定的 $1/\bar\epsilon$ 记号整理 $1/\epsilon-\gamma_E+\ln4\pi$，即得到\eqnrefs{eq:heat-kernel-divergence-master-dp11}。
\end{derivation}



% 四次标量相互作用的相关函数层级：已迁入 sections/01-exact-correlators-spectral.tex 的结论框及独立证明。














\begin{derivation}[从\eqnrefs{eq:phase2-si-start-r9}到\eqnrefs{eq:phase2-final-si-r9}]
正文\eqnrefs{eq:phase2-si-start-r9} 仍含两个末态三动量积分和一个四维守恒 $\delta$ 函数。在质心系中，三维 $\delta$ 函数先固定 $\mathbf p_4=-\mathbf p_3$，剩下的能量 $\delta$ 函数再固定径向动量模；最后只留下末态方向 $\dd\Omega$，从而得到正文\eqnrefs{eq:phase2-final-si-r9}。

采用能量四矢量
\begin{equation*}
 \mathcal P^\mu\equiv cp^\mu=(E,c\mathbf p),
 \end{equation*}
它只是物理四动量乘以常数 $c$。定义标准二体相空间因子
\begin{equation*}
 \dd\mathfrak P_2
 =(2\pi)^4\delta^{(4)}(\mathcal P_1+\mathcal P_2-\mathcal P_3-\mathcal P_4)
 \prod_{a=3,4}\frac{c^3\dd^3 p_a}{(2\pi)^3\,2E_a}.
 \end{equation*}
每个 $c^3\dd^3 p$ 都是 $\dd^3(c\mathbf p)$，所以 \eqnrefs{eq:phase2-si-start-r9} 完全由能量量纲的四动量组成。在质心系 $\mathcal P_1+\mathcal P_2=(c\sqrt s,\mathbf0)$。三维 $\delta$ 先令 $\mathbf p_4=-\mathbf p_3$，于是
\begin{equation*}
 \dd\mathfrak P_2
 =\frac{c^4p^2\dd p\,\dd\Omega}{(2\pi)^2 4E_3E_4}
 \delta(c\sqrt s-E_3-E_4).
\end{equation*}
令
\begin{equation*}
 f(p)=c\sqrt s-E_3(p)-E_4(p).
\end{equation*}
由 $\dd E_a/\dd p=c^2p/E_a$，
\begin{equation*}
 \left|f'(p)\right|
 =c^2p\left(\frac1{E_3}+\frac1{E_4}\right)
 =\frac{c^3p\sqrt s}{E_3E_4}.
\end{equation*}
所以
\begin{equation*}
 \delta[f(p)]
 =\frac{E_3E_4}{c^3p_f\sqrt s}\delta(p-p_f).
\end{equation*}
代回并积分 $p$，得到
\begin{equation*}
 {
 \dd\mathfrak P_2
 =\frac1{16\pi^2}\frac{|\mathbf p_f|}{\sqrt s}\,\dd\Omega.}
 \end{equation*}
右端无量纲，因为 $|\mathbf p_f|$ 与 $\sqrt s$ 都具有动量量纲。
\end{derivation}

\begin{derivation}[从\eqnrefs{eq:invariant-flux-scalar-r9}到\eqnrefs{eq:moller-velocity-si-r9}]
对两条入射轨迹，普通同一参考系中的相对相遇速率不能简单写成任意参考系都适用的 $|\mathbf v_1-\mathbf v_2|$。由四动量构造 Lorentz 标量

\begin{equation*}
 \mathcal I_{12}
 \equiv\sqrt{(p_1\!\cdot p_2)^2-m_1^2m_2^2c^4}.
 \end{equation*}
它具有动量平方量纲。定义 M\o ller 相对速率
\begin{equation*}
 v_{\rm M}
 \equiv
 \frac{c^3\mathcal I_{12}}{E_1E_2}.
 \end{equation*}
检查质心系时不跳过中间平方。写
\begin{equation*}
 p_1^\mu=\left(\frac{E_1}{c},\mathbf p_i\right),
 \qquad
 p_2^\mu=\left(\frac{E_2}{c},-\mathbf p_i\right),
\end{equation*}
于是
\begin{equation*}
 p_1\!\cdot p_2=\frac{E_1E_2}{c^2}+|\mathbf p_i|^2.
\end{equation*}
两条质量壳又分别给
\begin{equation*}
 m_a^2c^2=\frac{E_a^2}{c^2}-|\mathbf p_i|^2,
 \qquad a=1,2.
\end{equation*}
因此 \eqnrefs{eq:invariant-flux-scalar-r9} 根号中的量逐项化为
\begin{align*}
 \mathcal I_{12}^2
 &=\left(\frac{E_1E_2}{c^2}+p_i^2\right)^2
 -\left(\frac{E_1^2}{c^2}-p_i^2\right)
  \left(\frac{E_2^2}{c^2}-p_i^2\right),\\
 &=\frac{E_1^2E_2^2}{c^4}
 +\frac{2E_1E_2p_i^2}{c^2}+p_i^4
 -\frac{E_1^2E_2^2}{c^4}
 +\frac{(E_1^2+E_2^2)p_i^2}{c^2}-p_i^4,\\
 &=\frac{p_i^2}{c^2}(E_1+E_2)^2,
\end{align*}
其中 $p_i\equiv|\mathbf p_i|\ge0$。取正根得到
\begin{equation*}
 {
 \mathcal I_{12}
 =\frac{|\mathbf p_i|(E_1+E_2)}{c}
 =|\mathbf p_i|\sqrt s.}
 \end{equation*}
最后一步用了 $E_1+E_2=c\sqrt s$。于是
\begin{align*}
 v_{\rm M}
 &=|\mathbf p_i|c^2
 \left(\frac1{E_1}+\frac1{E_2}\right),\\
 &=|\mathbf v_1-\mathbf v_2|,
\end{align*}
因为 $\mathbf v_a=c^2\mathbf p_a/E_a$。所以 \eqnrefs{eq:moller-velocity-si-r9} 的协变表达确实回到熟悉的相对速度。
\end{derivation}

\begin{derivation}[从\eqnrefs{eq:delta-square-time-r15}到\eqnrefs{eq:2to2-cross-section-master-si-r9}]

取空间周期盒体积 $V$ 与观测时间 $T$。空间连续态满足
\begin{equation*}
 (2\pi\hbar)^3\delta^{(3)}(\mathbf0)=V,
 \qquad
 \sum_{\mathbf p}\longrightarrow
 V\int\frac{\dd^3 p}{(2\pi\hbar)^3}.
\end{equation*}
时间方向令 $\delta E\equiv E_f-E_i$，则
\begin{equation*}
 \int_{-T/2}^{T/2}\dd t\,\ee^{\ii\delta E t/\hbar}
 =T\,\operatorname{sinc}\!\left(\frac{\delta E T}{2\hbar}\right),
\end{equation*}
并由正文\eqnrefs{eq:delta-square-time-r15} 有
\begin{equation*}
 \lim_{T\to\infty}\frac1T
 \left|\int_{-T/2}^{T/2}\dd t\,\ee^{\ii\delta E t/\hbar}\right|^2
 =2\pi\hbar\delta(\delta E).
\end{equation*}
因此四维动量守恒 $\delta$ 函数平方在概率除以 $T$ 后留下一个能量守恒 $\delta$，其空间零点因子给出一个 $V$。结合两条初态的协变归一化，单位盒中单对入射粒子的二体末态跃迁率为正文\eqnrefs{eq:box-transition-rate-r15}，即
\begin{equation*}
 \dd\Gamma_{1+2\to2}
 =\frac{\hbar^2c^3}{4VE_1E_2}
 \overline{|\mathcal M|^2}\,\dd\mathfrak P_2.
\end{equation*}

单个归一化入射粒子在盒中的数密度是 $1/V$。正文\eqnrefs{eq:moller-velocity-si-r9} 给出相对速率，所以入射通量密度为
\begin{equation*}
 j_{12}=\frac{v_{\rm M}}{V}
 =\frac{c^3\mathcal I_{12}}{VE_1E_2}.
\end{equation*}
按正文中“跃迁率除以入射通量密度”的截面关系：
\begin{align*}
 \dd\sigma
 &=\frac{\dd\Gamma_{1+2\to2}}{j_{12}}\\
 &=\frac{\hbar^2}{4\mathcal I_{12}}
 \overline{|\mathcal M|^2}\,\dd\mathfrak P_2.
\end{align*}
$V$、$T$ 与参考系依赖的 $E_1E_2$ 在这一步全部消去。对 $2\to2$ 过程代入正文\eqnrefs{eq:phase2-final-si-r9}，并在质心系使用
$\mathcal I_{12}=|\mathbf p_i|\sqrt s$，得到
\begin{align*}
 \frac{\dd\sigma}{\dd\Omega}
 &=\frac{\hbar^2}{4|\mathbf p_i|\sqrt s}
 \overline{|\mathcal M|^2}
 \frac1{16\pi^2}\frac{|\mathbf p_f|}{\sqrt s}\\
 &=\frac{\hbar^2}{64\pi^2s}
 \frac{|\mathbf p_f|}{|\mathbf p_i|}
 \overline{|\mathcal M|^2},
\end{align*}
即正文\eqnrefs{eq:2to2-cross-section-master-si-r9}。
\end{derivation}

\begin{derivation}[从\eqnrefs{eq:delta-square-time-r15}到\eqnrefs{eq:master-decay-rate-si-r9}]
协变态归一化的单粒子初态在盒归一化后带因子 $(2P^0V)^{-1/2}$； 末态每条外腿给对应的质量壳测度。 LSZ 矩阵元中的四维 $\delta$ 在平方后，空间部分利用

\begin{equation*}
 (2\pi\hbar)^3\delta^{(3)}(\mathbf0)=V
\end{equation*}
消去母粒子的盒体积，而时间部分利用 \eqnrefs{eq:delta-square-time-r15} 把 $|\delta(E_f-E_i)|^2/T$ 化为单个能量 $\delta$。于是概率除以 $T$ 后只剩一个母粒子归一化因子 $1/(2P^0)$ 与 $n$ 体末态相空间。

为避免混淆能量量纲和动量量纲，先定义能量归一化截肢振幅
\begin{equation*}
 \mathcal A_E\equiv c\,\mathcal M.
\end{equation*}
若 $\mathcal M$ 在某个过程具有一个动量量纲，则 $\mathcal A_E$ 具有能量量纲。有限盒推导给
\begin{equation*}
 \dd\Gamma_{1\to n}
 =\frac{1}{2E_P\hbar}\,
 \overline{|\mathcal A_E|^2}\,\dd\mathfrak P_n,
 \qquad E_P=Mc^2.
 \end{equation*}
把 $\mathcal A_E=c\mathcal M$ 与 $E_P=Mc^2$ 代入，两个 $c^2$ 消去：
\begin{align*}
 \dd\Gamma_{1\to n}
 &=\frac{c^2}{2Mc^2\hbar}
 \overline{|\mathcal M|^2}\,\dd\mathfrak P_n,\\
 &={
 \frac{1}{2M\hbar}
 \overline{|\mathcal M|^2}\,\dd\mathfrak P_n.}
 \end{align*}
这里 $\dd\mathfrak P_n$ 使用能量四矢量 $\mathcal P^\mu=cp^\mu$ 的定义。以无量纲 Yukawa 耦合的费米子对衰变为例，振幅含有双线性因子 $\bar u v$；在上述外腿归一化下该双线性具有动量量纲，因此
\begin{equation*}
 \left[\frac{|\mathcal M|^2}{M\hbar}\right]
 =\mathrm{s^{-1}},
\end{equation*}
与衰变率的 SI 单位一致。这样 \eqnrefs{eq:master-decay-rate-si-r9} 与散射截面主式共享同一套有限盒、有限时间和外腿归一化，而不是独立记忆的另一条规则。
\end{derivation}



\begin{derivation}[从\eqnrefs{eq:phase3-recursive-dp52}到\eqnrefs{eq:phase3-dalitz-dp52}]
在母粒子静止系中，对正文\eqnrefs{eq:phase3-recursive-dp52} 的两个二体因子分别使用\eqnrefs{eq:phase2-final-si-r9}：

\begin{equation*}
\dd\mathfrak P_3
=\frac{\dd\mathcal S_{12}}{2\pi}
\frac{1}{16\pi^2}\frac{|\mathbf p_3|}{Mc}\dd\Omega_3
\frac{1}{16\pi^2}\frac{|\mathbf p_2^*|}{\sqrt{\mathcal S_{12}}/c}\dd\Omega_2^*.
\end{equation*}
若母粒子无自旋或已做自旋平均，整体方向积分给 $\int\dd\Omega_3=4\pi$。在 $q_{12}$ 静止系中以粒子 3 为极轴，方位角再给 $2\pi$，于是
\begin{equation*}
\dd\mathfrak P_3
=\frac{\dd\mathcal S_{12}\,\dd\cos\theta^*}{64\pi^3}
\frac{|\mathbf p_3|}{Mc}
\frac{|\mathbf p_2^*|}{\sqrt{\mathcal S_{12}}/c}.
\end{equation*}
正文\eqnrefs{eq:dalitz-angular-map-dp68} 给出
\begin{equation*}
\left|\frac{\partial\mathcal S_{23}}{\partial\cos\theta^*}\right|
=2c^2|\mathbf p_2^*||\mathbf p_3^*|.
\end{equation*}
因此把 $\cos\theta^*$ 换成 $\mathcal S_{23}$ 时，两个动量模恰好与 Jacobian 相消；再用 $\sqrt{\mathcal S_{12}}$ 和母粒子质量整理剩余因子，得到
\begin{equation*}
\dd\mathfrak P_3
=\frac{\dd\mathcal S_{12}\,\dd\mathcal S_{23}}{128\pi^3M^2c^4},
\end{equation*}
即正文\eqnrefs{eq:phase3-dalitz-dp52}。
\end{derivation}










